\section{Evaluation}\label{sec:evaluation}

\subsection{Test environment}

The evaluation runs on a four-node libvirt cluster of arm64
virtual machines (one master, three compute nodes), each booted
from a Yocto Scarthgap-built squashfs read-only root filesystem.
Each compute VM mounts a real \texttt{/dev/vdb} partition formatted
as FTRFS via virtio block. The host is a workstation running
Gentoo Linux with QEMU/TCG arm64 emulation (no hardware
passthrough): 20 host threads available, 31~GiB RAM, single SSD
backing all VM images. The VMs run linux-mainline 7.0.1 with the
hardening configuration documented in
\cref{sec:repro}; the ftrfs.ko module is identical
across all four nodes, verified by md5sum.

\subsection{Build hardening and CVE posture}

The build uses linux-mainline 7.0.1 (kernel.org stable point
release, SRCREV
\texttt{3cb1fb7a56d2fd8011f5282bc170c0d23dc1f4b5}) within Yocto
Scarthgap LTS. The \texttt{cve\_check} task reports zero
unpatched kernel CVEs and 12{,}562 patched after applying an
explicit exclusion file marking 88 pre-v7.0 NVD entries as
\texttt{fixed-version}; the exclusion file is committed at
\texttt{recipes-kernel/linux/cve-exclusions-kernel.inc} (commit
\texttt{9a5593b} of the
\texttt{yocto-hardened} layer~\cite{yocto-hardened-repo}). The
full image manifest reports 16{,}974 patched, 239 unpatched
(entirely user-space, scope of v0.4 per \cref{sec:future})
and 56 ignored, totalling 17{,}269 issues processed.

\subsection{Static and functional verification}

\begin{itemize}
\item \texttt{checkpatch.pl} with the kernel-tree configuration:
  zero errors and zero warnings on all modified filesystem source
  files at HEAD \texttt{9a63468}.
\item Compile-time sentinels (\texttt{BUILD\_BUG\_ON} in the
  kernel module, \texttt{static\_assert} in \texttt{mkfs.ftrfs})
  enforce on-disk struct sizes (24 bytes for
  \texttt{ftrfs\_rs\_event} in v3, 4096 bytes for
  \texttt{ftrfs\_super\_block}).
\item \texttt{generic/002}, \texttt{generic/010},
  \texttt{generic/098}, \texttt{generic/257} of the
  \texttt{xfstests} suite pass under manual invocation (4 of 5
  attempted; \texttt{generic/001} requires more than 2~GiB
  scratch, deferred).
\item Runtime \texttt{dmesg} sweep across all four nodes after a
  full bench pass: zero
  \texttt{BUG}/\texttt{WARN}/\texttt{Oops}/uncorrectable events.
\end{itemize}

\subsection{I/O baseline}

\Cref{tab:iobench} summarizes the I/O baseline measured on
2026-04-27 against the hardened build (squashfs MD5
\texttt{c56c203feb524f556329746fd693e805}). Each metric is sampled
30~times (3 compute nodes $\times$ 10~runs).

\begin{table*}[!t]
\centering
\small
\caption{FTRFS I/O baseline. 30 samples per metric. M1, M2 in
MB/s; M4 in seconds; M5 in ms/file. Bench JSON committed at
\texttt{papers/2026-04-ftrfs-v1/data/}.}
\label{tab:iobench}
\begin{tabular}{l l r r r r}
\toprule
\textbf{ID} & \textbf{Metric} &
\textbf{Min} & \textbf{Median} & \textbf{Max} & \textbf{Stddev} \\
\midrule
M1 & Write seq + fsync (4~MB)             & 2.128  & 5.000  & 5.556  & 0.584  \\
M2 & Read seq cold (4~MB)                 & 12.500 & 20.000 & 25.000 & 2.458  \\
M4 & Stat bulk (100 files)                & 0.140  & 0.150  & 0.160  & 0.006  \\
M5 & Small write + fsync (10$\times$64~B) & 21.000 & 23.000 & 30.000 & 2.331  \\
\bottomrule
\end{tabular}
\end{table*}

The metric distributions are concentrated, reflecting the
deterministic behaviour of the metadata paths under the load
profile (per axis~3.4). The discrete clustering of M1 and M2
values is a measurement artifact of \texttt{dd}'s temporal
resolution at the millisecond granularity. The single M1 outlier
at 2.128~MB/s corresponds to a host load spike during run capture
(load average 1.99/2.45/3.44), included in the data without
correction in the interest of full disclosure. A bench harness
with sub-millisecond timing is on the roadmap for v2.

Throughput is not the design objective: a fair comparison against
\texttt{ext4 + dm-integrity + dm-verity} on the same platform is
reserved for v3, when stage~4 enables feature parity.

Slurm cluster scheduling latencies, captured in the same session,
are reported in \cref{tab:slurm}; they document the
deployment context end-to-end.

\begin{table}[ht]
\centering
\small
\caption{Slurm scheduler latency on the FTRFS-backed cluster.}
\label{tab:slurm}
\begin{tabular}{l l r}
\toprule
\textbf{ID} & \textbf{Phase} & \textbf{Latency (s)} \\
\midrule
S1 & Job submit latency (1 node)        & 0.300 \\
S2 & Parallel job (3 nodes)             & 0.370 \\
S3 & Throughput (9 jobs)                & 5.020 \\
S4 & FTRFS write from Slurm job         & PASS \\
\bottomrule
\end{tabular}
\end{table}
